\chapter{Urinary Incontinence}

\section*{Definition}
Involuntary leakage urine affecting quality life stress urge mixed predominant types.

\section*{Pathophysiology}
Differentiate stress leakage cough sneeze urge sudden strong desire leakage overflow retention prostate neurologic etiations varying causes require management strategies lifestyle modifications therapeutic interventions surgical options.

\section*{Etiology}
\begin{itemize}
  \item Stress weakness women prostate enlargement
  \item Overactive bladder detrusor overactivity
  \item Neurogenic bladder spinal cord injury stroke Parkinson MS
  \item Overflow benign prostatic hyperplasia neurogenic
  \item Functional mobility limitation cognitive impairment environmental barriers
\end{itemize}

\section*{Indications for Evaluation}
Seek care if:
\begin{itemize}
  \item Severe or worsening symptoms
  \item Symptoms lasting more than Evaluate causing bother embarrassment recurrent UTIs skin breakdown social isolation psychological distress
  \item Accompanied by other concerning signs
\end{itemize}

\section*{Related Symptoms}
\begin{itemize}
  \item Frequency urgency nocturia dysuria recurrent infections bedwetting dribbling reduced stream hesitancy post-void residual sensation
\end{itemize}
\textbf{Note:} Always consider serious underlying conditions when evaluating persistent urinary incontinence.

\section*{Self-Care / Home Management}
\begin{itemize}
  \item Drink adequate water throughout the day unless fluid restriction is medically indicated.
  \item Avoid bladder irritants including caffeine, alcohol, acidic foods, and artificial sweeteners.
  \item Practice good hygiene and urinate promptly when the urge arises to prevent urinary stasis.
  \item Apply a heating pad to the lower abdomen or back for comfort during urinary discomfort.
  \item Seek medical attention for fever, flank pain, visible blood in urine, or difficulty urinating.
\end{itemize}

\section*{Diagnostic Approach}
\begin{itemize}
  \item History covers urinary symptoms (frequency, urgency, dysuria, hesitancy, stream changes), fever, and flank pain.
  \item Urinalysis with microscopy and culture is the first-line diagnostic test for urinary symptoms.
  \item Imaging (renal ultrasound, CT urogram) evaluates for stones, obstruction, structural abnormalities, or masses.
  \item Urodynamic studies or cystoscopy may be indicated for complex or refractory voiding dysfunction.
\end{itemize}

\section*{Treatment / Management Overview}
\begin{itemize}
  \item Urinary tract infections are treated with targeted antibiotics based on culture sensitivity results.
  \item Nephrolithiasis management includes hydration, analgesics, medical expulsive therapy, or urologic intervention for large stones.
  \item Voiding dysfunction may respond to behavioral modification, pelvic floor therapy, or pharmacotherapy.
  \item Urology referral is indicated for recurrent infections, hematuria, obstructive symptoms, or suspected malignancy.
\end{itemize}

\clearpage
